\documentclass[border=10pt]{standalone}
\usepackage{circuitikz}

\begin{document}
\color{white}
\begin{circuitikz}[scale=2, transform shape, line width=1pt, american voltages]

% Op-amp node
    \node[op amp] (op) at (4.5,0) {};

    % Non-inverting (+) input connected to ground
    \draw (op.+) node[ground] {};

    % Junction coordinate where the two input branches meet
    \coordinate (in_junction) at (2.4, 0.5);

    % Top input branch: V1 -> R1 -> junction
    \draw (0, 0.9) node[left] {$V_1$} to[short, o-] ++(0.3,0)
          to[R, l=$R_1$] ++(1.7,0) |- (in_junction);

    % Bottom input branch: V2 -> R2 -> junction
    \draw (0, 0.1) node[left] {$V_2$} to[short, o-] ++(0.3,0)
          to[R, l=$R_2$] ++(1.7,0) |- (in_junction);

    % Connection from junction to Node X and inverting (-) input
    \draw (in_junction) -- (3.2, 0.5) node[circ] (X) {} node[above left] {$X$}
          -- (op.-);

    % Output wire with junction dot for feedback connection
    \draw (op.out) -- ++(0.5,0) node[circ] (out_node) {}
          to[short, -o] ++(0.6,0) node[right] {$V_{out}$};

    % Feedback path with resistor RF aligned using orthogonal coordinates
    \draw (X) -- (X |- 0, 1.5)
          to[R, l=$R_F$] (out_node |- 0, 1.5)
          -- (out_node);

\end{circuitikz}
\end{document}
